\appendix


\clearpage

\section{Additional Experimental Setups}
\label{appendix:setup}
\begin{figure}
    \centering
    \includegraphics[width=0.95\linewidth]{figure/perf_time_xl.pdf}
    \vspace{-0.1in}
    \caption{QA performance (F1) and efficiency (Time/Query) for different retrieval-augmented generation approaches. We use the FLAN-T5-XL (3B) as the base LLM.}
    \label{fig:perf_effi_xl}
    % \vspace{-0.175in}
\end{figure}


\subsection{Datasets}
We use publicly open datasets for both single-hop and multi-hop QA datasets, referring to as~\citet{DPR} and \citet{ircot}, respectively. We describe the characteristics of each dataset:

\noindent
\textbf{1) SQuAD v1.1}~\cite{squad} is created through a process where annotators write questions based on the documents they read. 

\noindent
\textbf{2) Natural Questions}~\cite{nq} is constructed by real user queries on Google Search. 

\noindent
\textbf{3) TriviaQA}~\cite{tqa} comprises trivia questions sourced from various quiz websites.

\noindent
\textbf{4) MuSiQue}~\cite{musique} is collected by compositing multiple single-hop queries, to form queries spanning 2-4 hops. 

\noindent
\textbf{5) HotpotQA}~\cite{hotpotqa} is constructed by having annotators create questions that link multiple Wikipedia articles. 

\noindent
\textbf{6) 2WikiMultiHopQA}~\cite{2hop} is derived from Wikipedia and its associated knowledge graph path, needing 2-hops.

% For the test data, we use the subsampled data from the development sets.


\subsection{Models}
We describe the details of models as follows:

\noindent
\textbf{1) No Retrieval.} This approach uses only the LLM itself, to generate the answer to the given query. 

\noindent
\textbf{2) Single-step Approach.} This approach first retrieves the relevant knowledge with the given query from the external knowledge sources and then augments the LLM with this retrieved knowledge to generate the answer, which iterates only once.

\noindent
\textbf{3) Adaptive Retrieval.} This baseline~\cite{adaptiveRetrieval} adaptively augments the LLM with the retrieval module, only when the entities appearing in queries are less popular. To extract entities, we use the available entity-linking method~\cite{blinke}, namely BLINK, for questions.

\noindent
\textbf{4) Self-RAG.} This baseline~\cite{self-rag} trains the LLM to adaptively perform retrieval and generation, where the retrieval is conducted once it predicts the special retrieval token above a certain threshold, and the answer generation follows.

\noindent
\textbf{5) Adaptive-RAG.} This is our model that adaptively selects the retrieval-augmented generation strategy, smoothly oscillating between the non-retrieval, single-step approach, and multi-step approaches\footnote{For the multi-step approach, we use the state-of-the-art question answering strategy from IRCoT~\cite{ircot}.} without architectural changes, based on the query complexity assessed by the classifier. 

\noindent
\textbf{6) Multi-step Approach.} This approach~\cite{ircot} is the multi-step retrieval-augmented LLM, which iteratively accesses both the retriever and LLM with interleaved Chain-of-Thought reasoning~\cite{cot} repeatedly until it derives the solution or reaches the maximum step number.

\noindent
\textbf{7) Adaptive-RAG w/ Oracle} This is an ideal scenario of our Adaptive-RAG equipped with an oracle classifier that perfectly categorizes the query complexity. 


\subsection{Implementation Details}
For computing resources, we use A100 GPUs with 80GB memory. In addition, due to the significant costs associated with evaluating retrieval-augmented generation models, we perform experiments with a single run. Finally, we implemented models using PyTorch~\cite{DBLP:conf/nips/PaszkeGMLBCKLGA19} and Transformers library~\cite{DBLP:conf/emnlp/WolfDSCDMCRLFDS20}.

\section{Additional Experimental Results}

\begin{figure}
    \centering
    \includegraphics[width=0.95\linewidth]{figure/perf_time_xxl.pdf}
    \vspace{-0.1in}
    \caption{QA performance (F1) and efficiency (Time/Query) for different retrieval-augmented generation approaches. We use the FLAN-T5-XXL (11B) as the base LLM.}
    \label{fig:perf_effi_xxl}
    % \vspace{-0.175in}
\end{figure}



\begin{table*}[t!]
\caption{Results on each of a collection of datasets with FLAN-T5-XXL (11B) as the LLM. We emphasize our results in bold.}
\vspace{-0.1in}
\label{tab:main:xxl}
\small
\centering
\resizebox{\textwidth}{!}{
\renewcommand{\arraystretch}{1.0}
\begin{tabular}{lllccccccccccccccc}
\toprule

& & & \multicolumn{5}{c}{\bf SQuAD} & \multicolumn{5}{c}{\bf Natural Questions} & \multicolumn{5}{c}{\bf TriviaQA} \\
\cmidrule(l{2pt}r{2pt}){4-8} \cmidrule(l{2pt}r{2pt}){9-13} \cmidrule(l{2pt}r{2pt}){14-18}
\textbf{Data} & \textbf{Types} &\textbf{Methods} & EM & F1 & Acc & Step & Time & EM & F1 & Acc & Step & Time & EM & F1 & Acc & Step & Time \\

\midrule

\multirowcell{7}[-0.0ex][l]{\textbf{Single-step}} 

& \multirowcell{2}[-0.0ex][l]{\textbf{Simple}} 

& \textbf{No Retrieval} & 7.00 & 14.40 & 8.40 & 0.00 & 0.08 & 18.80 & 25.50 & 20.40 & 0.00 & 0.08 & 32.80 & 39.20 & 35.40 & 0.00 & 0.08 \\

& & \textbf{Single-step Approach} & 28.80 & 40.80 & 35.00 & 1.00 & 1.00 & 41.40 & 51.20 & 47.60 & 1.00 & 1.00 & 56.00 & 64.70 & 61.80 & 1.00 & 1.00 \\

\noalign{\vskip 0.25ex}\cdashline{2-18}\noalign{\vskip 0.75ex}

& \multirowcell{3}[-0.0ex][l]{\textbf{Adaptive}} 

& \textbf{Adaptive Retrieval} & 15.60 & 25.60 & 20.00 & 0.50 & 0.54 & 31.00 & 39.70 & 35.00 & 0.50 & 0.54 & 44.80 & 52.20 & 48.60 & 0.50 & 0.54 \\

& & \textbf{Self-RAG$^*$} & 1.60 & 11.90 & 20.80 & 0.59 & 0.31 & 39.20 & 47.10 & 42.40 & 0.75 & 0.09 & 14.60 & 33.70 & 60.20 & 0.76 & 0.22 \\

& & \textbf{Adaptive-RAG (Ours)} & \textbf{27.80} & \textbf{39.80} & \textbf{34.00} & \textbf{1.17} & \textbf{1.50} & \textbf{41.20} & \textbf{51.00} & \textbf{47.40} & \textbf{1.00} & \textbf{1.00} & \textbf{52.00} & \textbf{60.30} & \textbf{57.20} & \textbf{1.03} & \textbf{1.33} \\

\noalign{\vskip 0.25ex}\cdashline{2-18}\noalign{\vskip 0.75ex}

& \textbf{Complex}

& \textbf{Multi-step Approach} & 24.60 & 36.90 & 30.20 & 2.13 & 3.83 & 39.60 & 49.60 & 46.40 & 2.16 & 3.94 & 52.60 & 61.10 & 59.40 & 2.17 & 4.03 \\

\noalign{\vskip 0.25ex}\cdashline{2-18}\noalign{\vskip 0.75ex}

& \textbf{Oracle}
& \textbf{Adaptive-RAG w/ Oracle} & 32.80 & 46.90 & 38.20 & 0.85 & 0.94 & 51.20 & 61.00 & 57.00 & 0.71 & 0.91 & 63.40 & 71.30 & 68.20 & 0.51 & 0.60 \\

\midrule
\midrule

& & & \multicolumn{5}{c}{\bf MuSiQue} & \multicolumn{5}{c}{\bf HotpotQA} & \multicolumn{5}{c}{\bf 2WikiMultiHopQA} \\
\cmidrule(l{2pt}r{2pt}){4-8} \cmidrule(l{2pt}r{2pt}){9-13} \cmidrule(l{2pt}r{2pt}){14-18}
\textbf{Data} & \textbf{Types} &\textbf{Methods} & EM & F1 & Acc & Step & Time & EM & F1 & Acc & Step & Time & EM & F1 & Acc & Step & Time \\

\midrule

\multirowcell{7}[-0.0ex][l]{\textbf{Multi-step}} 

& \multirowcell{2}[-0.0ex][l]{\textbf{Simple}} 

& \textbf{No Retrieval} & 4.20 & 13.40 & 5.40 & 0.00 & 0.08 & 17.40 & 25.44 & 18.40 & 0.00 & 0.09 & 26.80 & 32.93 & 28.00 & 0.00 & 0.08 \\

& & \textbf{Single-step Approach} & 16.80 & 25.70 & 19.20 & 1.00 & 1.00 & 37.60 & 49.27 & 39.60 & 1.00 & 1.00 & 46.60 & 54.13 & 48.20 & 1.00 & 1.00 \\

\noalign{\vskip 0.25ex}\cdashline{2-18}\noalign{\vskip 0.75ex}

& \multirowcell{3}[-0.0ex][l]{\textbf{Adaptive}} 

& \textbf{Adaptive Retrieval} & 8.40 & 17.80 & 10.20 & 0.50 & 0.54 & 26.60 & 36.01 & 27.80 & 0.50 & 0.54 & 35.20 & 42.68 & 36.80 & 0.50 & 0.54 \\

& & \textbf{Self-RAG$^*$} & 1.20 & 8.20 & 11.80 & 0.68 & 0.27 & 5.60 & 17.86 & 30.60 & 0.76 & 0.26 & 3.00 & 19.14 & 39.00 & 0.90 & 0.25 \\

& & \textbf{Adaptive-RAG (Ours)} & \textbf{20.60} & \textbf{28.50} & \textbf{23.20} & \textbf{1.89} & \textbf{3.12} & \textbf{44.20} & \textbf{54.78} & \textbf{46.80} & \textbf{1.58} & \textbf{2.53} & \textbf{47.60} & \textbf{57.36} & \textbf{54.00} & \textbf{1.46} & \textbf{2.55} \\

\noalign{\vskip 0.25ex}\cdashline{2-18}\noalign{\vskip 0.75ex}

& \textbf{Complex}

& \textbf{Multi-step Approach} & 19.40 & 27.50 & 21.80 & 2.09 & 3.66 & 47.00 & 57.81 & 49.40 & 2.08 & 3.73 & 57.60 & 67.65 & 64.00 & 2.17 & 3.63 \\

\noalign{\vskip 0.25ex}\cdashline{2-18}\noalign{\vskip 0.75ex}

& \textbf{Oracle}
& \textbf{Adaptive-RAG w/ Oracle} & 24.20 & 37.20 & 26.60 & 1.22 & 1.71 & 52.20 & 64.80 & 54.60 & 0.92 & 1.33 & 59.20 & 70.40 & 68.60 & 0.82 & 1.14 \\

\bottomrule

\end{tabular}
}
% \vspace{-0.125in}
\end{table*}


\begin{table*}[t!]
\caption{Results on each of a collection of datasets with GPT-3.5 (Turbo) as the LLM. We emphasize our results in bold.}
\vspace{-0.1in}
\label{tab:main:gpt}
\small
\centering
\resizebox{\textwidth}{!}{
\renewcommand{\arraystretch}{1.0}
\begin{tabular}{lllccccccccccccccc}
\toprule

& & & \multicolumn{5}{c}{\bf SQuAD} & \multicolumn{5}{c}{\bf Natural Questions} & \multicolumn{5}{c}{\bf TriviaQA} \\
\cmidrule(l{2pt}r{2pt}){4-8} \cmidrule(l{2pt}r{2pt}){9-13} \cmidrule(l{2pt}r{2pt}){14-18}
\textbf{Data} & \textbf{Types} &\textbf{Methods} & EM & F1 & Acc & Step & Time & EM & F1 & Acc & Step & Time & EM & F1 & Acc & Step & Time \\

\midrule

\multirowcell{7}[-0.0ex][l]{\textbf{Single-step}} 

& \multirowcell{2}[-0.0ex][l]{\textbf{Simple}} 

& \textbf{No Retrieval} & 16.00 & 29.20 & 23.80 & 0.00 & 0.62 & 39.80 & 55.70 & 55.00 & 0.00 & 0.56 & 64.00 & 75.60 & 75.80 & 0.00 & 0.68 \\

& & \textbf{Single-step Approach} & 18.00 & 33.80 & 29.20 & 1.00 & 1.00 & 32.40 & 46.80 & 54.80 & 1.00 & 1.00 & 55.20 & 66.50 & 65.80 & 1.00 & 1.00 \\

\noalign{\vskip 0.25ex}\cdashline{2-18}\noalign{\vskip 0.75ex}

& \multirowcell{3}[-0.0ex][l]{\textbf{Adaptive}} 

& \textbf{Adaptive Retrieval} & 15.40 & 30.00 & 24.40 & 0.50 & 0.81 & 36.40 & 51.20 & 56.60 & 0.50 & 0.78 & 62.00 & 71.90 & 72.20 & 0.50 & 0.84 \\

& & \textbf{Self-RAG$^*$} & 1.60 & 11.90 & 20.80 & 0.59 & 1.91 & 39.20 & 47.10 & 42.40 & 0.75 & 0.52 & 14.60 & 33.70 & 60.20 & 0.76 & 1.59 \\

& & \textbf{Adaptive-RAG (Ours)} & \textbf{19.80} & \textbf{34.40} & \textbf{30.00} & \textbf{0.87} & \textbf{1.21} & \textbf{36.80} & \textbf{52.00} & \textbf{56.60} & \textbf{0.68} & \textbf{0.86} & \textbf{62.40} & \textbf{73.80} & \textbf{73.80} & \textbf{0.22} & \textbf{0.79} \\

\noalign{\vskip 0.25ex}\cdashline{2-18}\noalign{\vskip 0.75ex}

& \textbf{Complex}

& \textbf{Multi-step Approach} & 17.40 & 31.50 & 26.20 & 2.50 & 3.24 & 35.60 & 49.70 & 57.80 & 2.58 & 3.79 & 54.80 & 67.10 & 68.00 & 2.30 & 2.65 \\

\noalign{\vskip 0.25ex}\cdashline{2-18}\noalign{\vskip 0.75ex}

& \textbf{Oracle}
& \textbf{Adaptive-RAG w/ Oracle} & 28.00 & 45.90 & 39.40 & 0.54 & 0.93 & 50.00 & 65.40 & 67.00 & 0.28 & 0.8 & 70.80 & 81.00 & 80.00 & 0.11 & 0.73 \\

\midrule
\midrule

& & & \multicolumn{5}{c}{\bf MuSiQue} & \multicolumn{5}{c}{\bf HotpotQA} & \multicolumn{5}{c}{\bf 2WikiMultiHopQA} \\
\cmidrule(l{2pt}r{2pt}){4-8} \cmidrule(l{2pt}r{2pt}){9-13} \cmidrule(l{2pt}r{2pt}){14-18}
\textbf{Data} & \textbf{Types} &\textbf{Methods} & EM & F1 & Acc & Step & Time & EM & F1 & Acc & Step & Time & EM & F1 & Acc & Step & Time \\

\midrule

\multirowcell{7}[-0.0ex][l]{\textbf{Multi-step}} 

& \multirowcell{2}[-0.0ex][l]{\textbf{Simple}} 

& \textbf{No Retrieval} & 20.40 & 31.30 & 24.40 & 0.00 & 0.81 & 37.40 & 51.04 & 43.20 & 0.00 & 0.74 & 37.00 & 48.50 & 43.40 & 0.00 & 0.90 \\

& & \textbf{Single-step Approach} & 16.40 & 26.70 & 23.60 & 1.00 & 1.00 & 39.60 & 50.44 & 45.60 & 1.00 & 1.00 & 46.80 & 57.69 & 52.60 & 1.00 & 1.00 \\

\noalign{\vskip 0.25ex}\cdashline{2-18}\noalign{\vskip 0.75ex}

& \multirowcell{3}[-0.0ex][l]{\textbf{Adaptive}} 

& \textbf{Adaptive Retrieval} & 18.80 & 30.30 & 24.80 & 0.50 & 0.90 & 38.60 & 50.70 & 43.20 & 0.50 & 0.87 & 44.20 & 55.11 & 50.60 & 0.50 & 0.95 \\

& & \textbf{Self-RAG$^*$} & 1.20 & 8.20 & 11.80 & 0.68 & 1.66 & 5.60 & 17.86 & 30.60 & 0.76 & 1.67 & 3.00 & 19.14 & 39.00 & 0.90 & 1.81 \\

& & \textbf{Adaptive-RAG (Ours)} & \textbf{21.80} & \textbf{32.60} & \textbf{29.60} & \textbf{1.90} & \textbf{2.29} & \textbf{40.40} & \textbf{52.56} & \textbf{47.00} & \textbf{0.93} & \textbf{1.48} & \textbf{46.60} & \textbf{60.09} & \textbf{56.80} & \textbf{1.59} & \textbf{2.23} \\

\noalign{\vskip 0.25ex}\cdashline{2-18}\noalign{\vskip 0.75ex}

& \textbf{Complex}

& \textbf{Multi-step Approach} & 23.00 & 32.50 & 31.60 & 3.41 & 3.61 & 45.80 & 58.36 & 52.20 & 2.73 & 3.18 & 52.20 & 66.08 & 62.40 & 3.36 & 3.35 \\

\noalign{\vskip 0.25ex}\cdashline{2-18}\noalign{\vskip 0.75ex}

& \textbf{Oracle}
& \textbf{Adaptive-RAG w/ Oracle} & 29.60 & 44.70 & 35.60 & 0.90 & 1.45 & 55.60 & 69.90 & 62.80 & 0.54 & 1.08 & 52.20 & 69.90 & 66.60 & 0.65 & 1.21 \\

\bottomrule

\end{tabular}
}
% \vspace{-0.125in}
\end{table*}





\paragraph{Performance vs Time}
We further provide a comparison of different retrieval-augmented generation approaches with FLAN-T5-XL and FLAN-T5-XXL models in Figure~\ref{fig:perf_effi_xl} and Figure~\ref{fig:perf_effi_xxl}, respectively, in the context of performance and efficiency trade-offs. Similar to the observation made from the GPT-3.5 model in Figure~\ref{fig:perf_effi_gpt}, our proposed Adaptive-RAG is significantly more effective as well as efficient.

\paragraph{Performance per Dataset}
In addition to detailing the performance of each dataset with the FLAN-T5-XL model, as shown in Table \ref{tab:main:xl}, we also present the results for each dataset with the FLAN-T5-XXL and GPT-3.5 models in Table~\ref{tab:main:xl} and Table~\ref{tab:main:gpt}, respectively. The experimental results show that our Adaptive-RAG consistently balances between efficiency and accuracy. It is worth noting that while the GPT-3.5 model performs effectively in addressing straightforward queries even without document retrieval, it benefits significantly from our Adaptive-RAG in terms of effectiveness when solving complex multi-hop queries.



